\section{Writing It Down}
\label{sec:ch18-writing-it-down}

\index{trace id!written down everywhere|(}%
Section~\ref{sec:ch18-the-id-was-already-there} ends with an id that reaches
every process and is recorded by one of them. The fix adds no tracing. It is three small edits that make the id visible at
the three places it currently disappears: the response, the services' own log output, and the
statement log this book counts with.

\subsection{The Router Hands It Back}

\index{config.yaml@\code{config.yaml}!returning the trace id|(}%
\code{response\_trace\_id}, under \code{telemetry.tracing}, is the router's
setting for the first one. It puts the id the router is already using into a
response header, named \code{x-wg-trace-id} unless you name it otherwise~\autocite{cosmotelemetry}; HTTP header names are
case-insensitive, and the router writes this one as \code{X-Wg-Trace-Id}. Here is \code{router/config.yaml}, complete,
with the block this chapter adds:

\begin{minted}[highlightlines={50-59}]{yaml}
version: "1"

# Without this the router asks Cosmo Cloud for its execution config and
# refuses to start, because it has no token to ask with. The path is
# resolved against the directory the router is started in, not against
# this file, so start it from this folder:
#
#   router -config config.yaml
execution_config:
  file:
    path: ../graph/router.json

# A development switch, and what makes the query plan visible at all.
# The plan is an Advanced Request Tracing option, and outside dev mode
# the router wants a Cosmo Cloud token before it hands one out. With
# this off, X-WG-Include-Query-Plan is ignored rather than refused.
dev_mode: true

# Who the router thinks is asking. The alternative entry in this list
# is a jwks url, which is what a deployment with an identity provider
# behind it uses; the shared secret is here because this book stands
# up no such provider. header_key_id is not optional: the router
# matches it against the kid in the token header, so the token this
# book mints carries one.
#
# A request carrying a token this key does not verify is answered 401
# by the router and never reaches a subgraph.
authentication:
  jwt:
    jwks:
      - secret: splitting-the-graph-development-signing-key-32b
        symmetric_algorithm: HS256
        header_key_id: dev

# The router does not forward request headers to subgraphs unless it
# is told to, and a subgraph that guards its own data cannot tell who
# is asking without this block. Leaving it out does not make the graph
# insecure; it makes the guarded field unreachable. Every request that
# names one comes back null with an authorization error underneath,
# for a caller holding a good token exactly as for a caller holding
# none, which is what makes the missing block read like a policy
# rather than like a header. A request that never names a guarded
# field is unaffected, so this survives a smoke test.
headers:
  all:
    request:
      - op: propagate
        named: Authorization

# The router gives every request a trace id, and sends it to each
# subgraph it fetches from in a traceparent header, with nothing
# configured here at all. What it does not do is tell the client.
# This puts the id on every response as X-WG-Trace-Id, which is the
# one string that joins a report from a client to the router's log
# and to the log of every service that request reached.
telemetry:
  tracing:
    response_trace_id:
      enabled: true
\end{minted}

\index{X-WG-Trace-Id@\code{X-WG-Trace-Id}}%
Nothing else under \code{telemetry} needs to be set for it, and the page from
section~\ref{sec:ch18-the-page-that-got-slower} now comes back with one more
header than it did:

\begin{minted}{text}
X-Wg-Trace-Id: 327f111c797fd261f05d6bce80d680ca
\end{minted}

That is the string a client can put in a bug report, and it is the value the
router writes on its own line for the request.
\index{config.yaml@\code{config.yaml}!returning the trace id|)}%

\subsection{The Services Print Their Scopes}

\index{appsettings.json@\code{appsettings.json}!printing log scopes|(}%
The second place is every log entry a service writes through .NET's logger,
which already carries the id in its scope. The console only has to be told to
print it. That is configuration rather than code, and it goes in
\code{appsettings.json}, the settings file \code{dotnet new web} generates,
which every service in this book has carried since it was created and none
has needed until now. This is the whole file, and it is identical in all four services:

\begin{minted}{json}
{
  "Logging": {
    "LogLevel": {
      "Default": "Information",
      "Microsoft.AspNetCore": "Warning"
    },
    "Console": {
      "FormatterName": "simple",
      "FormatterOptions": {
        "IncludeScopes": true
      }
    }
  },
  "AllowedHosts": "*"
}
\end{minted}

\index{logging!a formatter option that is ignored}%
Both settings under \code{Console} are needed.

\code{FormatterOptions} is where the console formatter's settings live, and
written on its own, without \code{FormatterName}, it is read by nothing: no
error, no warning, and no scopes. With the formatter named, every entry the
service logs gains a line carrying \code{TraceId}. None of this book's services
logs much at the levels that file lets through, so the place those scopes earn their keep is
section~\ref{sec:ch18-when-it-fails-instead}, when something fails.
\index{appsettings.json@\code{appsettings.json}!printing log scopes|)}%

\subsection{The Statement Log Prints the Id}

\index{StatementLog@\code{StatementLog}!prints the trace id|(}%
The third place is the instrument this book counts with, and it needs a code
change because it writes to the console directly. Here is the Sessions
service's \code{StatementLog.cs}, complete, with what changed marked: it reads
the trace id off the current activity and prints it on the marker line.

\begin{minted}[highlightlines={2,11-13,42,43}]{csharp}
using System.Data.Common;
using System.Diagnostics;
using Microsoft.EntityFrameworkCore.Diagnostics;

namespace Sessions;

/// <summary>
/// Prints every statement the service sends to SQLite, numbered, so a
/// request's cost can be read off the console. The counter is
/// process-wide and never reset: start the service, send one request,
/// and the last number is what that request cost. Each marker carries
/// the trace id of the request that issued it, so the statements one
/// request cost can be picked out of a log that many requests share.
/// </summary>
public sealed class StatementLog : DbCommandInterceptor
{
    private static int _count;

    public override InterceptionResult<DbDataReader> ReaderExecuting(
        DbCommand command,
        CommandEventData eventData,
        InterceptionResult<DbDataReader> result)
    {
        Print(command);
        return result;
    }

    public override ValueTask<InterceptionResult<DbDataReader>>
        ReaderExecutingAsync(
            DbCommand command,
            CommandEventData eventData,
            InterceptionResult<DbDataReader> result,
            CancellationToken cancellationToken = default)
    {
        Print(command);
        return ValueTask.FromResult(result);
    }

    private static void Print(DbCommand command)
    {
        var n = Interlocked.Increment(ref _count);
        var trace = Activity.Current?.TraceId.ToString() ?? "none";
        Console.WriteLine($"-- statement {n} trace {trace}");
        Console.WriteLine(command.CommandText);
    }
}
\end{minted}

\index{StatementLog@\code{StatementLog}!four identical copies}%
The Speakers and Ratings copies are this file with their own namespace, as
they have been since chapter~\ref{ch:second-third-subgraph}, and so now is the
Search service's, which chapter~\ref{ch:cross-seam-paging} wrote in a shorter
shape that this edit replaces. The prefix \code{-- statement} is unchanged, so
counting markers still counts statements. A statement issued outside any
request prints the word \code{none} where the id would be.
\index{StatementLog@\code{StatementLog}!prints the trace id|)}%

\subsection{Finding the Service}

\index{blame routing!finding the slow service by its id|(}%
Now run the investigation from
section~\ref{sec:ch18-the-page-that-got-slower} again, on the same four
services with chapter~\ref{ch:router-n-plus-1}'s naive resolver still in
Speakers and these three edits made. The client reports the page as slow and
quotes the header above. Search each of the four consoles for that id, and
count the statement markers carrying it:

\begin{center}
\begin{tabular}{lr}
\toprule
Service & Statements carrying the id \\
\midrule
Sessions, on 5001 & 1 \\
Speakers, on 5002 & 3 \\
Ratings, on 5003 & 1 \\
Search, on 5004 & 0 \\
\bottomrule
\end{tabular}
\end{center}

\index{N+1@N+1!found by a trace id}%
The count is the same whether that console printed one request or a thousand
around it, because the id selects the request's statements rather than
whatever happened to be adjacent. Here are the three from Speakers, marker
lines only, the query under each being the one-row \code{SELECT}
section~\ref{sec:ch18-the-page-that-got-slower} printed:

\begin{minted}{text}
-- statement 4 trace 327f111c797fd261f05d6bce80d680ca
-- statement 5 trace 327f111c797fd261f05d6bce80d680ca
-- statement 6 trace 327f111c797fd261f05d6bce80d680ca
\end{minted}

\index{ownership!the investigation ends with a team}%
Speakers is the service, the Speakers team is who gets the ticket, and the
numbers \code{4} to \code{6} show the other half of the argument: this console
had already printed three statements for an earlier request, and without the
id nothing would separate those three from these. What the Speakers team does
next is chapter~\ref{ch:router-n-plus-1}'s section on the loader behind the
resolver. What this chapter changes is that they are the team who reads it,
rather than whoever owns the router.
\index{blame routing!finding the slow service by its id|)}%
\index{trace id!written down everywhere|)}%
